\chapter{隣接多項間漸化式}\label{章：隣接多項間漸化式}
\begin{bookcolumns}
隣接3項間漸化式は,隣り合う二項を一組にすれば連立漸化式になる.
本章ではこの対応から特性方程式を導き,通常の解法に現れる補助数列の意味を説明する.
その後,差分によって項数を減らす方法を調べ,一般の項数へ広げよう.
\section{隣接3項間漸化式}\label{節：隣接3項間漸化式}
\subsection{隣り合う二項を一組にする}

前章の連立漸化式に対し,今度は一つの数列について
\[
 a_{n+2}=5a_{n+1}-6a_n
\]
を考える.
$a_n$一つだけでは次の状態を決められないが,\,隣り合う二項を持てばよい.
\[
 \boldsymbol v_n=\begin{pmatrix}a_{n+1}\\a_n\end{pmatrix},
 \qquad B=\begin{pmatrix}5&-6\\1&0\end{pmatrix}
\]
とすると
\[
 B\boldsymbol v_n
 =\begin{pmatrix}5a_{n+1}-6a_n\\a_{n+1}\end{pmatrix}
 =\boldsymbol v_{n+1}.
\]
第二行は,\,上の成分を一つ古い記録として下へ移す.
隣り合う二項から次を決める式も,\,\textbf{二つの成分を同時に更新する式}である.

前章と同じように,行列を掛けると一定倍率になるベクトルを探す.
この例では
\[
 B\begin{pmatrix}2\\1\end{pmatrix}
 =2\begin{pmatrix}2\\1\end{pmatrix},
 \qquad
 B\begin{pmatrix}3\\1\end{pmatrix}
 =3\begin{pmatrix}3\\1\end{pmatrix}.
\]
二方向に分ければ,\,一方は2倍,\,他方は3倍されるだけになる.

\subsection{特性方程式は倍率を決める}\label{節：特性方程式は倍率を決める}

一般の隣接3項間漸化式
\[
 a_{n+2}=s a_{n+1}+t a_n
\]
では,$B=\begin{pmatrix}s&t\\1&0\end{pmatrix}$である.
固有ベクトルを$\begin{pmatrix}x\\y\end{pmatrix}$とすると
\[
 sx+ty=\lambda x,\qquad x=\lambda y.
\]
$y=0$なら$x=0$にもなるので,\,固有ベクトルでは$y\ne0$である.
$y$で割って$\begin{pmatrix}\lambda\\1\end{pmatrix}$にそろえると,
\[
 s\lambda+t=\lambda^2
\]
が必要になる. 従って,\,倍率は
\[
 \lambda^2-s\lambda-t=0
\]
という\textbf{特性方程式}から決まる.
その根を\textbf{特性根}とも呼ぶ.
これが,\,漸化式の解法に二次方程式が現れる理由である.

\subsection{二つの等比数列に分ける}

二つの異なる根を$\lambda_1,\lambda_2$とする.
対応する$\begin{pmatrix}\lambda_1\\1\end{pmatrix}$と
$\begin{pmatrix}\lambda_2\\1\end{pmatrix}$は比例しないので,\,初期状態を一意に
\[
 \begin{pmatrix}a_2\\a_1\end{pmatrix}
 =c\begin{pmatrix}\lambda_1\\1\end{pmatrix}
 +d\begin{pmatrix}\lambda_2\\1\end{pmatrix}
\]
と分けられる.
各方向を$n-1$回更新して下の成分を読むと,
\[
 a_n=c\lambda_1^{n-1}+d\lambda_2^{n-1}.
\]
根が非零なら通常の冪として読み,\,根が零ならその成分は初期状態にだけ寄与し,\,一回の更新で消えると読む.
初期値を合わせた後は更新の一意性によって,\,全項がこの式に一致する.

\noindent\textbf{【例題】}\par\nobreak
$a_1=1,\,a_2=4$,$a_{n+2}=5a_{n+1}-6a_n$の一般項を求めよ.

\noindent\textbf{【方針】}\par\nobreak
2倍と3倍の方向に初期状態を分けます.

\noindent\textbf{【解答】}\par\nobreak
\[
 a_n=-2^{n-1}+2\cdot3^{n-1}.
\]
\noindent\textbf{【解法】}\par\nobreak
特性根は2,3である. 初期状態は
\[
 \begin{pmatrix}4\\1\end{pmatrix}
 =-\begin{pmatrix}2\\1\end{pmatrix}
 +2\begin{pmatrix}3\\1\end{pmatrix}.
\]
従って
\[
 \boldsymbol v_n
 =-2^{n-1}\begin{pmatrix}2\\1\end{pmatrix}
 +2\cdot3^{n-1}\begin{pmatrix}3\\1\end{pmatrix}.
\]
下の成分から一般項を得る.
$n=1,2$で初期値に合い,\,二つの冪はそれぞれ更新式を満たす.
線形性によって,\,その和も更新式を満たす.

\subsection{補助数列を取り出す}

同じ式は
\[
 a_{n+2}-2a_{n+1}
 =3(a_{n+1}-2a_n)
\]
と変形しても解ける.
この$a_{n+1}-2a_n$は,\,行列のどこにあるのだろうか.

数を横に並べた$L=(u,v)$を\textbf{行ベクトル}と呼び,
$L\begin{pmatrix}x\\y\end{pmatrix}=ux+vy$と定める.
先ほどの$B$について
\[
 (1,-2)B=(3,-6)=3(1,-2)
\]
だから
\[
 \begin{aligned}
 b_n&=(1,-2)\boldsymbol v_n=a_{n+1}-2a_n,\\
 b_{n+1}&=(1,-2)B\boldsymbol v_n=3b_n.
 \end{aligned}
\]
このように$LA=\lambda L$を満たす零でない行ベクトルを\textbf{左固有ベクトル}という.

固有ベクトルは,\,一定倍率になる\textbf{状態の方向}を見つける.
左固有ベクトルは,\,一定倍率になる\textbf{取り出し方}を見つける.
二つの解法は,\,同じ更新を別の側から見ている.

\subsection{重解で現れる足し算}\label{節：重解で現れる足し算}

特性方程式が重解$r\ne0$を持つと
\[
 a_{n+2}=2r a_{n+1}-r^2a_n.
\]
$a_n=r^{n-1}b_n$と置いて倍率を取り除けば
\[
 b_{n+2}-2b_{n+1}+b_n=0.
\]
これは$b_{n+2}-b_{n+1}=b_{n+1}-b_n$である.
差が一定なので
\[
 b_n=C+D(n-1),
\]
従って
\[
 a_n=r^{n-1}\{C+D(n-1)\}.
\]
等比的な変化を取り除いた後に,\,等差的な変化が残った.
重解だから同じ冪を二回書くのでは,\,二つの初期値を自由に合わせられない.

行列でも,\,この足し算を直接見られる.
例えば$r=2$では
\[
 B=\begin{pmatrix}4&-4\\1&0\end{pmatrix},
 \quad
 \boldsymbol u=\begin{pmatrix}2\\1\end{pmatrix},
 \quad
 \boldsymbol w=\begin{pmatrix}1\\0\end{pmatrix}
\]
とすると
\[
 B\boldsymbol u=2\boldsymbol u,\qquad
 B\boldsymbol w=\boldsymbol u+2\boldsymbol w.
\]
$\boldsymbol w$は固有ベクトルではないが,\,更新すると$\boldsymbol u$の方向へも寄与する.
このように$(B-rI)\boldsymbol w=\boldsymbol u$,
$(B-rI)\boldsymbol u=\boldsymbol0$となる$\boldsymbol w$は,\,\textbf{一般化固有ベクトル}の例である.

$\boldsymbol v_n=c_n\boldsymbol u+d_n\boldsymbol w$と書くと
\[
 d_{n+1}=2d_n,\qquad c_{n+1}=2c_n+d_n.
\]
各係数から$2^{n-1}$を取り除けば,\,第二の式は一定量を足す式になる.
毎回の寄与が積み重なるため,$n$が現れるのである.

\noindent\textbf{【例題】}\par\nobreak
$a_1=1,\,a_2=4$,$a_{n+2}=4a_{n+1}-4a_n$の一般項を求めよ.

\noindent\textbf{【方針】}\par\nobreak
公比2の成長を取り除き,\,差を調べます.

\noindent\textbf{【解答】}\par\nobreak
\[
 a_n=n2^{n-1}.
\]
\noindent\textbf{【解法】}\par\nobreak
$b_n=\frac{a_n}{2^{n-1}}$とすると
$b_1=1,\,b_2=2$で,\,二階差が零になる.
従って$b_n=n$である.
行列では$\boldsymbol v_1=\boldsymbol u+2\boldsymbol w$に対応する.

なお,\,一般の行列は固有値が重なっても,\,固有ベクトルを十分に選べる場合がある.
例えば$2I$はどの方向も2倍にする.
この場合は足し算の蓄積がなく,\,すべての成分が単に2倍される.
\textbf{重解そのものと,\,固有方向が足りないことは区別する}.
$r=0$の隣接3項間漸化式は$a_{n+2}=0$なので,\,最初の二項だけを別に記録すればよい.

\subsection{二つの固有方向が得られない場合}

対角化できるために必要なのは,\,状態を一意に表せるだけの固有ベクトルである.
例えば
\[
 A=\begin{pmatrix}r&1\\0&r\end{pmatrix}
\]
では,\,固有値は$r$だけである.
$(A-rI)\begin{pmatrix}x\\y\end{pmatrix}=\boldsymbol0$から$y=0$となるため,\,固有方向は一本しかない.
二成分の状態全体を固有ベクトルだけでは表せない.

しかし,\,この表示のままでも
\[
 y_{n+1}=ry_n,\qquad x_{n+1}=rx_n+y_n
\]
と順に解ける.
$r\ne0$では
\[
 \begin{aligned}
 y_n&=r^{n-1}y_1,\\
 x_n&=r^{n-1}x_1+(n-1)r^{n-2}y_1.
 \end{aligned}
\]
これは直前に見た重解の場合と同じ蓄積である.
$r=0$では二回の更新で両成分が零になる.

\begin{bookpoint}
対角化できれば各成分が独立する.
できない場合にも,\,ある成分から別の成分への寄与を順に追えることがある.
等比的な倍率に重なる多項式は,\,その寄与の蓄積を表す.
\end{bookpoint}

\section{差分と項数の削減}\label{節：差分と項数の削減}
\subsection{差を取ると項数が減る理由}

前節では$a_{n+1}-2a_n$を一つの数列として追った.
まず倍率が1の場合,つまり隣り合う二項の差を考えよう.
初期値$a_1=1,a_2=2,a_3=6$の下で
\[
 a_{n+3}=6a_{n+2}-11a_{n+1}+6a_n
\]
を左辺へまとめると
\[
 \begin{aligned}
 &(a_{n+3}-a_{n+2})-5(a_{n+2}-a_{n+1})\\
 &\qquad+6(a_{n+1}-a_n)=0.
 \end{aligned}
\]
従って$b_n=a_{n+1}-a_n$と置けば
\[
 b_{n+2}-5b_{n+1}+6b_n=0
\]
という隣接3項間漸化式になる.
元の隣接4項間漸化式より,必要な初期値が一つ少ない.
ただし差だけでは元の数列を一意に決められない.
$b_n$を解いた後に
\[
 a_n=a_1+\sum_{j=1}^{n-1}b_j
\]
と戻すため,残る初期値$a_1$も使う.
$n=1$では和を0とする.


初期値から$b_1=1,b_2=4$となる.
前節の例題と同じ式・初期値なので,
\[
 b_n=-2^{n-1}+2\cdot3^{n-1}.
\]
和を取ると
\[
 \begin{aligned}
 a_n&=1-(2^{n-1}-1)+(3^{n-1}-1)\\
 &=1-2^{n-1}+3^{n-1}.
 \end{aligned}
\]
差分を取る前の初期値$a_1$を用いて,定数の成分も取り戻せた.

なぜこの例では差が有効なのだろうか.
対応する多項式は
\[
 X^3-6X^2+11X-6=(X-1)(X-2)(X-3)
\]
であり,差を取ることが因子$X-1$に対応する.

\subsection{項を一つ進める操作}

数列の添字を一つ進める操作を$E$と書く.
すなわち$(Ea)_n=a_{n+1}$であり,
\[
 ((E-r)a)_n=a_{n+1}-ra_n
\]
と定める. ここで$E$は数ではなく,数列に作用する記号である.
二度進める操作は$E^2$,何もしない操作は$I$と書く.
表記を簡単にするため,$E-rI$を$E-r$と略している.

例えば前の隣接4項間漸化式は
\[
 (E^3-6E^2+11E-6)a=0
\]
であり,多項式と同じように因数分解して
\[
 (E-2)(E-3)\{(E-1)a\}=0
\]
と書ける. 右辺の0は全項が0の数列を表す.
$b=(E-1)a$と置けば,残った因子が$b$の漸化式を与える.
各因子は同じ$E$の多項式なので,作用させる順序を交換できる.

一般に,特性多項式$P(X)$が
\[
 P(X)=(X-r)Q(X)
\]
と分かれるなら,$b_n=a_{n+1}-ra_n$によって
\[
 P(E)a=0\quad\Longleftrightarrow\quad Q(E)b=0
\]
となる. 特に通常の差分$\Delta a_n=a_{n+1}-a_n$が使えるのは,$P(1)=0$の場合である.
根が1でない式に,通常の差分を取るだけで項数が減るとは限らない.

\subsection{倍率を取り除いてから差を取る}

$r\ne0$とし,$a_n=r^{n-1}c_n$と置けば
\[
 a_{n+1}-ra_n=r^n(c_{n+1}-c_n).
\]
従って$a_{n+1}-ra_n$も,倍率$r$の等比的な変化を取り除いた後の差分である.
異なる根を一つずつ取り除く方法も,重解の場合に差を繰り返す方法も,同じ因子の計算で説明できる.

\begin{bookpoint}
行列では,更新を一定倍率の成分に分ける.
差分では,その倍率に対応する因子を一つずつ取り除く.
項数が減る代わりに,元へ戻すときに初期値を一つ使う.
\end{bookpoint}

\section{一般の隣接多項間漸化式}\label{節：一般の隣接多項間漸化式}
\subsection{過去を状態にまとめる}

次の項に過去三項が必要なら,\,その三項を状態として持つ.
例えば
\[
 a_{n+3}=6a_{n+2}-11a_{n+1}+6a_n
\]
は
\[
 \begin{pmatrix}a_{n+3}\\a_{n+2}\\a_{n+1}\end{pmatrix}
 =
 \begin{pmatrix}6&-11&6\\1&0&0\\0&1&0\end{pmatrix}
 \begin{pmatrix}a_{n+2}\\a_{n+1}\\a_n\end{pmatrix}.
\]
第一行で新しい項を作り,\,下の行で記録を一つずつ移す.
二成分の場合から増えたのは,\,記憶する過去の長さである.

一般に定数係数の隣接$k+1$項間線形漸化式
\[
 a_{n+k}=\sum_{j=0}^{k-1}c_j a_{n+j}
\]
に対して
\[
 \boldsymbol v_n=
 \begin{pmatrix}a_{n+k-1}\\a_{n+k-2}\\\vdots\\a_n\end{pmatrix}
\]
と置くと,$\boldsymbol v_{n+1}=C\boldsymbol v_n$である.
更新行列は
\[
 C=
 \begin{pmatrix}
 c_{k-1}&c_{k-2}&\cdots&c_1&c_0\\
 1&0&\cdots&0&0\\
 0&1&\cdots&0&0\\
 \vdots&&\ddots&&\vdots\\
 0&0&\cdots&1&0
 \end{pmatrix}.
\]
成分数が増えても,\,ベクトルの和や行列の積は各行ごとに同じ方法で定める.
この形の行列を\textbf{コンパニオン行列}という.

\subsection{特性方程式の一般化}

固有ベクトルの一番下の成分を1にそろえると,\,下の行から順に
\[
 \boldsymbol u(\lambda)=
 \begin{pmatrix}\lambda^{k-1}\\\lambda^{k-2}\\\vdots\\1\end{pmatrix}
\]
という形になる.
一番下が零なら上も順に零となるため,\,このそろえ方ができる.
第一行の条件から
\[
 \lambda^k=\sum_{j=0}^{k-1}c_j\lambda^j
\]
を得る. 従って特性多項式は
\[
 \chi(X)=X^k-\sum_{j=0}^{k-1}c_jX^j.
\]
項数が増えても,\,特性方程式$\chi(\lambda)=0$は一定倍率の方向を探している.
高次方程式の根を簡単な式で求められるかどうかと,\,この構造が成り立つかどうかは別の問題である.

\noindent\textbf{【例題】}\par\nobreak
$a_1=1,\,a_2=2,\,a_3=6$とし,
\[
 a_{n+3}=6a_{n+2}-11a_{n+1}+6a_n
\]
の一般項を求めよ.

\noindent\textbf{【方針】}\par\nobreak
前節では差分で解いた例を,今度は固有値から直接解きます. 特性多項式を因数分解し,\,三つの等比的な成分に分けます.

\noindent\textbf{【解答】}\par\nobreak
\[
 a_n=1-2^{n-1}+3^{n-1}.
\]
\noindent\textbf{【解法】}\par\nobreak
\[
 X^3-6X^2+11X-6=(X-1)(X-2)(X-3).
\]
従って$a_n=A+B2^{n-1}+D3^{n-1}$と置く.
初めの三項から
\[
 A+B+D=1,\quad A+2B+3D=2,
\]
\[
 A+4B+9D=6.
\]
これを解くと$A=1,\,B=-1,\,D=1$となる.
得た式は初期三項と更新を満たすので,\,全項に一致する.

\subsection{初期条件と解の自由度}\label{節：初期条件と解の自由度}

$k+1$個の隣接項を結ぶ更新では,\,初めの$k$項を自由に選ぶと,\,その後は順に一意に決まる.
従って解の数列全体は,\,初期値の組$\mathbb C^k$と一対一に対応する.
二つの解の和と定数倍も解であるため,\,初期値だけでなく解全体にも線形性がある.
ここで「解全体は$k$次元」というのは,\,\textbf{独立に選べる初期の数が$k$個ある}という意味である.

特性根$\lambda_1,\ldots,\lambda_k$が相異なり非零なら,
\[
 a_n=\sum_{j=1}^{k}A_j\lambda_j^{n-1}
\]
で全解を表せる.
なぜこれで十分なのかを確かめよう.
各冪は更新を満たす. また,\,異なる冪の数列は一次独立である.

その独立性は,前節で定義した操作$E$を使うと示しやすい.
\[
 (E-\mu)\lambda^{n-1}
 =(\lambda-\mu)\lambda^{n-1}.
\]
もし$\sum_j A_j\lambda_j^{n-1}$が零数列なら,\,固定した$j$以外の根について
$E-\lambda_\ell$をすべて作用させる.
$j$以外の項が消え,
\[
 A_j\prod_{\ell\ne j}(\lambda_j-\lambda_\ell)
 \lambda_j^{n-1}=0
\]
となる. $n=1$を見れば$A_j=0$である.
根が異なるので積は零でない.

従って$k$本は独立で,\,自由度$k$の解全体の基底になる.
初期値を合わせる係数も一意に決まる.
一般項に必要な任意定数の個数は,\,記憶すべき初期条件の個数と対応する.

\subsection{重根と繰り返し取った差分}

非零の根$r$が重なると,\,冪$r^{n-1}$だけでは自由度が不足する.
そこで,\,多項式$P(n)$を掛けた
\[
 a_n=r^{n-1}P(n)
\]
を調べる. $\Delta P(n)=P(n+1)-P(n)$を\textbf{前進差分}と書くと,
\[
 (E-r)a_n=r^n\Delta P(n).
\]
もう一度作用させれば$r^{n+1}\Delta^2P(n)$になる.
一般に
\[
 (E-r)^m a_n=r^{n+m-1}\Delta^mP(n).
\]
差分を一回取ると非定数多項式の次数が一つ下がるので,\,次数$m-1$以下なら$\Delta^mP=0$である.
従って,$r$が$m$重根なら
\[
 r^{n-1},\ (n-1)r^{n-1},\ \ldots,\
 (n-1)^{m-1}r^{n-1}
\]
が解になる.

すべての根が非零で,\,異なる根を$r_1,\ldots,r_h$,
その重複度を$m_1,\ldots,m_h$とすれば,
\[
 a_n=\sum_{j=1}^{h}r_j^{n-1}P_j(n),
 \qquad \deg P_j<m_j
\]
で全解を表せる.
ここで$\sum_jm_j=k$であり,\,係数の総数も$k$である.
これらの独立性も差分で示せる.
ある根$r_j$に注目し,\,他の根の因子$(E-r_\ell)^{m_\ell}$を作用させると他の成分は消える.
一方,$r_j$の多項式成分では,\,異なる根の因子は最高次係数を非零倍するため次数を保つ.
残った成分が零数列ならその多項式は零であり,\,各根について同じことがいえる.
独立な解が$k$本あるので全解を尽くす.

例えば
\[
 a_{n+3}-3a_{n+2}+3a_{n+1}-a_n=0
\]
では,\,特性多項式は$(X-1)^3$である.
三階差が零なので
\[
 a_n=A+B(n-1)+D(n-1)(n-2)
\]
となる. 重根1に対応する多項式は,\,足し算を重ねた数列でもある.

零が$m_0$重根の場合には,\,最初の$m_0$項だけ自由でその後が零になる数列も解になる.
特性多項式を$\chi(X)=X^{m_0}Q(X)$とすると,
$Q(E)$と$E^{m_0}$は交換でき,\,そのような数列は$E^{m_0}a=0$を満たすためである.
非零根の成分にこれらを加えれば,\,初期条件の自由度は合計$k$になる.
零の冪や負の指数へ無理に代入しないことが大切である.

\subsection{非同次の式と一般化の範囲}

これまでのように項の線形結合だけで書ける式を\textbf{同次}という.
そこに外から指定された付加項がある式を\textbf{非同次}という.
右辺に既知の数列$f_n$が加わる
\[
 a_{n+k}=\sum_{j=0}^{k-1}c_j a_{n+j}+f_n
\]
でも,\,一つの解$u_n$が分かれば,\,差$b_n=a_n-u_n$は
\[
 b_{n+k}=\sum_{j=0}^{k-1}c_j b_{n+j}
\]
を満たす.
従って解は,\,\textbf{一つの特解と同次の一般解の和}になる.
不動点を引く方法は,\,特解として定数列を選んだ場合だった.

例えば
\[
 a_{n+2}-3a_{n+1}+2a_n=1
\]
では$u_n=-n$が特解である.
従って$a_n=A+B2^{n-1}-n$と書ける.
更新の異なる部分を先に一つ引き受けてもらい,\,残りを分解するのである.

なお,\,係数が$n$によって変わる場合も
$\boldsymbol v_{n+1}=C_n\boldsymbol v_n$と記録できるが,
\[
 \boldsymbol v_n=C_{n-1}\cdots C_2C_1\boldsymbol v_1
\]
となる. 毎回の行列を個別に対角化しても,\,同じ座標が次の更新に適するとは限らない.
本節の一定倍率への分解が,\,そのまま使えるわけではない.

\begin{bookpoint}
一つの数列の隣接項を追う式も,複数の数列を追う式も,成分をまとめれば行列で更新できる.
定数係数の線形漸化式では,\,状態を一つの行列で更新し,\,その倍率と成分間の寄与を追う.
\end{bookpoint}

\end{bookcolumns}
